\section{Q\&A: costos, funciones objetivo y pluralidad de valores}

Las slides preparadas terminan alrededor de 52:40; las preguntas posteriores vuelven una y otra vez a gráficas anteriores, pero añaden las restricciones más importantes de la clase. Esta sección conserva solo la nueva spoken evidence y no repite las figuras.

\subsection{Descriptive Morality y Prescription se separan de nuevo}

Un estudiante pregunta si Delphi toma los juicios existentes como la moral que debe seguirse. La ponente responde que el objetivo de la investigación es, ante todo, modelar descriptive judgments, y reconoce que estos contienen bias; un normative endorsement requiere principios adicionales y procedimientos sociales. También describe ChatGPT como un generation model más potente, no como un knowledge model ya terminado.

\teachervoice{La aclaración central del Q\&A es esta: poder predecir «qué diría la gente» no equivale a saber «qué debe hacerse». Aunque la prediction accuracy sea alta, el sistema aún debe mostrar de qué grupos provienen los datos, cómo se distribuyen los desacuerdos y en qué condiciones se abstiene de responder.}

\subsection{Culture, Dataset y Distribution of Opinions}

Los datos morales provienen de distintas épocas, países, plataformas y annotation instructions. La ponente señala que Delphi tiene un Western bias y menciona también que los distintos moral datasets difieren en escala y cobertura. Un objetivo más razonable no es emitir una respuesta única para todas las cuestiones en disputa, sino distinguir entre unacceptable harm y legitimate disagreement, y representar la opinion distribution.

\begin{knowledgebox}{Representación del Pluralism en aprendizaje automático}
Una classification de etiqueta única reduce el desacuerdo a un voto mayoritario. Una salida más completa puede incluir label distribution, condiciones por grupo, confidence, rationale clusters y un contested flag. Esto sigue sin resolver los conflictos normativos, pero al menos no oculta el disagreement bajo una apariencia de certeza.
\end{knowledgebox}

\subsection{Complexity: Maieutic e Hybrid son lentos}

La explanation expansion recursiva llama al LM muchas veces, el weighted MaxSAT debe resolverse sobre un grafo y el hybrid añade además retrieval y principle checking. La ponente admite abiertamente que tanto Maieutic prompting como Delphi hybrid son lentos. Por ello, la validación del sistema debe reportar token calls, tree width/depth, solver time y end-to-end latency, y no solo accuracy.

\subsection{La comparación entre COMET y GPT-3 no es una Architecture Ablation}

El Q\&A confirma de nuevo que COMET usa supervisión específica de ATOMIC, que GPT-3 usa un few-shot prompt y que el critical teacher pasa además por critic filtering. Los resultados pueden respaldar que «la especialización en la tarea y la calidad de los datos tienen valor», pero no permiten estimar por separado la contribución causal del parameter count, del pretraining corpus o de la architecture.

\subsection{Por qué los humanos también necesitan un Critic}

La ponente explica el critic a partir del proceso de aprendizaje: los humanos no creen para siempre lo primero que se les ocurre, sino que forman un internal critic mediante retroalimentación, reflexión y corrección social. Esta analogía aporta intuición de diseño, pero la distribución de entrenamiento del critic de una máquina aún debe auditarse de forma explícita; no adquiere credibilidad por el solo hecho de «parecerse a un humano».

\subsection{La Refusal no sustituye a la Context Understanding}

Un estudiante pregunta si el sistema debería evitar por completo las moral / ethical questions. La ponente considera que el moral advice directo puede restringirse, pero que el sistema aún debe reconocer contextos de acoso, peligro, odio o petición de ayuda; de lo contrario, depender solo de la abstención impide que los componentes de seguridad detecten el harmful context real. En ingeniería, esto exige separar la answer policy del understanding model.

\begin{importantbox}{La Abstention también es una capacidad que debe entrenarse y evaluarse}
Una buena abstención no consiste en detenerse al encontrar una palabra clave, sino en poder expresar la incertidumbre, reconocer el alto riesgo, ofrecer alternativas seguras, registrar la causa que la activó y escalar a una persona las preguntas que requieren un juicio autorizado. La abstention debe reportar dos tipos de error: false refusal y unsafe answer.
\end{importantbox}

\subsection{Next-Token, Preference y Knowledge son objetivos distintos}

Al final del Q\&A se vuelve a la tesis más profunda de esta clase: la next-token prediction aprende la continuación del lenguaje, el RLHF / preference tuning aprende las preferencias humanas y el knowledge model aprende relaciones del mundo e inferencias consultables. Los tres comparten representaciones y se refuerzan entre sí, pero no son el mismo objective. El commonsense puede emerge a partir de un LM, pero emerge no equivale a estable, calibrable ni con capacidad de auditoría.

\teachervoice{La ponente afirma que el espacio del adversarial commonsense es prácticamente infinito; ante una combinación absurda que escuchan por primera vez, los humanos pueden juzgarla gracias a su conceptual understanding, mientras que el Transformer tiende más a predecir qué palabra sigue. Esta es la razón última por la que sostiene que «un language model no equivale a un knowledge model».}

\subsection{Las Humanities no son un External Reviewer, sino codiseñadoras}

La última pregunta, sobre value pluralism, extiende la responsabilidad a las humanities at large. Qué diferencias constituyen un desacuerdo aceptable y qué daños deben rechazarse no puede decidirlo solo la comunidad de AI researchers. La ponente usa el derecho como analogía: las normas son un artifact humano formado por negociación social; los modelos pueden ayudar a representarlas y verificarlas, pero la legitimidad proviene de la participación institucional.

\subsection{Resumen de la sección}

El Q\&A vuelve a situar las propuestas técnicas de la prepared talk en el contexto de los recursos, los objetivos y la gobernanza. Maieutic / hybrid tienen un costo computacional considerable; la specialist and generalist comparison tiene factores de confusión; abstenerse no equivale a comprender; y la pluralidad de valores requiere representaciones distribucionales y participación institucional interdisciplinaria.
